\FloatBarrier
\subsection{Symbolic Axis Mapping}
\label{sec:axis-mapping}
Symbolic axis mapping connects groups defined by different operations. It
records where each group appears in a tensor, allowing composition of value
expressions while preserving axis order and operand pairings.

\FloatBarrier
\subsubsection{Building and Composing Mapping Tables}
\label{sec:mapping-tables}
Each record $(P,G,p)$ satisfies Equation~\eqref{eq:mapping}. We display records
in tables whose columns denote tensors, whose first row gives the current
interface, and whose other rows give group positions. A dash means that the
example has no mapping for that group. Tables can include both registered
and derivable mappings.

\begin{figuretext}{figures/Symbolic Axis Mapping.pdf}{Shared tensor interfaces connect local axis descriptions without changing tensor values or physical axis order.}{fig:axis-mapping}{Local dot and addition mapping tables are connected through their shared
intermediate, producing a table whose axes use a common coordinate description.}
Figure~\ref{fig:axis-mapping} first illustrates the principle with
\code{T = dot(X,Y);} \code{O = add(T,Z)}. The dot describes its output using $I$ and
$J$, while addition describes that same interface using $S$. Connecting
$S\leftrightarrow I\catseq J$ through $T$ aligns the coordinate descriptions
without changing the arithmetic expression or physical axis order. The
concrete matrix positions below the symbolic expressions are examples,
rather than fixed-rank restrictions on the method. Unlike the following
worked example, the figure places the dot inside the addition; both require
connecting local descriptions through a shared tensor.


We now work through the distribution example
\code{T = Add(X,Y);} \code{O = dot_general(T,Z)} from
Section~\ref{sec:composition}. Let $\rho(G)=\operatorname{range}(|G|)$ denote
all positions in $G$. Ranks and position attributes remain symbolic.
\paragraph{1. Build addition's interface.}
Addition binds \code{X} and \code{Y} to one ordered group $S$ and yields $S$.
Its reads and result thus have these mappings:
\begin{center}
\begin{tabular}{@{}lccc@{}}\toprule
Interface/group & $X$ & $Y$ & $T$\\\midrule
Axis interface & $S$ & $S$ & $S$\\
$S$ & $\rho(S)$ & $\rho(S)$ & $\rho(S)$\\\bottomrule
\end{tabular}
\end{center}

\end{figuretext}

\paragraph{2. Establish the dot's operand layouts.}
The dot introduces batch group $B$, contraction group $K$, and free groups
$I,J$. Let $b_L,b_R$ be its batch attributes, $c_L,c_R$ its contraction
attributes, and $R$ the initial interface of \code{Z}. Free positions are
\begin{equation}
\begin{aligned}
 f_L&=\ocomp(\rho(S),b_L\catseq c_L),\\
 f_R&=\ocomp(\rho(R),b_R\catseq c_R).
\end{aligned}
\end{equation}
Bindings require valid position lists and matching paired extents and must
cover each operand without overlap. Groups can be interleaved, so concatenating
$B,K,I$ does not necessarily reproduce the left operand's dimension order.
Instead, binding positions restore the original order:
\begin{equation}
\begin{aligned}
 \widehat S&=\gath(B\catseq K\catseq I,
                  \operatorname{inverse}(b_L\catseq c_L\catseq f_L)),\\
 \widehat R&=\gath(B\catseq K\catseq J,
                  \operatorname{inverse}(b_R\catseq c_R\catseq f_R)).
\end{aligned}
\label{eq:rebuilt-interfaces}
\end{equation}
Each concatenated position sequence is a permutation because the bindings
form a complete disjoint partition. Thus
\begin{equation}
\begin{aligned}
 \gath(\widehat S,b_L)&=B,&
 \gath(\widehat S,c_L)&=K,&
 \gath(\widehat S,f_L)&=I.
\end{aligned}
\end{equation}
The right interface satisfies analogous bindings. Registering these mappings
and partitions gives
\begin{center}
\begin{tabular}{@{}lcc@{}}\toprule
Interface/group & $T$ & $Z$\\\midrule
Axis interface & $\widehat S$ & $\widehat R$\\
$B$ & $b_L$ & $b_R$\\
$K$ & $c_L$ & $c_R$\\
$I$ & $f_L$ & --\\
$J$ & -- & $f_R$\\\bottomrule
\end{tabular}
\end{center}
The interfaces describe the operands in dot coordinates; they do not change
physical axis order. An existing producer interface is reused when it already
satisfies the bindings.

\paragraph{3. Apply the dot's mappings to \code{X} and \code{Y}.}
Addition describes \code{T} using $S$, whereas the dot describes those same
axes using $\widehat S$. Since addition preserves order, each axis of its
output corresponds to the same-position axes in \code{X} and \code{Y}.
Transporting its reads into $\widehat S$ therefore places $B,K,I$ at positions
$b_L,c_L,f_L$ in all three tensors. The compiler can now compose
\code{Reduce(K, Product(Add(X,Y), Z))}, with $K$ referring to corresponding
coordinates in \code{X}, \code{Y}, and \code{Z}.

\paragraph{4. Extend the table to the output.}
A later consumer sees the dot's output, where surviving axes occur in output
order and $K$ is absent. Registering the order $B\catseq I\catseq J$ extends
the table as follows. The row $\widehat S$ denotes addition's whole-axis group
after coordinate transport.
\begin{center}
\small
\setlength{\tabcolsep}{4pt}
\begin{tabular}{@{}lccccc@{}}\toprule
Interface/group & $X$ & $Y$ & $T$ & $Z$ & $O$\\\midrule
Axis interface & $\widehat S$ & $\widehat S$ & $\widehat S$ & $\widehat R$ & $B\catseq I\catseq J$\\
$\widehat S$ & $\rho(\widehat S)$ & $\rho(\widehat S)$ & $\rho(\widehat S)$ & -- & --\\
$B$ & $b_L$ & $b_L$ & $b_L$ & $b_R$ & $\rho(B)$\\
$K$ & $c_L$ & $c_L$ & $c_L$ & $c_R$ & --\\
$I$ & $f_L$ & $f_L$ & $f_L$ & -- & $|B|+\rho(I)$\\
$J$ & -- & -- & -- & $f_R$ & $|B|+|I|+\rho(J)$\\\bottomrule
\end{tabular}
\end{center}
Adding a length to a position sequence shifts each entry. The entries for
\code{X} and \code{Y} follow from transporting the dot partition into
addition's reads; output entries follow from the declared order.

\paragraph{5. Find positions through existing mappings.}
Not every relationship needs to be stored directly. Suppose $p_X$ locates
$\widehat S$ in \code{X}, and $c_L$ locates $K$ in $\widehat S$:
\begin{equation}
 \gath(P_X,p_X)=\widehat S,\qquad \gath(\widehat S,c_L)=K.
\end{equation}
Composing the selections yields
\begin{equation}
\begin{aligned}
 K&=\gath(\widehat S,c_L)\\
  &=\gath(\gath(P_X,p_X),c_L)\\
  &=\gath(P_X,\gath(p_X,c_L)).
\end{aligned}
\label{eq:gather-composition}
\end{equation}
Thus the $K,X$ entry is $\gath(p_X,c_L)$. For addition, $p_X$ lists all
positions in order, simplifying this expression to $c_L$. Repeated composition
recovers positions through longer chains. Addition's single group makes the
correspondence straightforward; selecting across several producer groups
requires partitioning them to expose finer relationships.

\FloatBarrier
\subsubsection{Symbolic Axis Partitioning}
\label{sec:partitioning}
Figure~\ref{fig:axis-partitioning} illustrates the information retained when a
consumer selects across producer groups. Three positions must be distinguished:
a selected axis's position within its producer group, its slot in the
consumer's contraction list, and its partner's position in the other operand.

\paperfigure{1}{figures/Symbolic Axis Partitioning (Dot Reassociation).pdf}{Partitioning selected axes by producer origin preserves their contraction partners for reassociation.}{fig:axis-partitioning}{Three panels show selecting producer output axes, splitting them by origin,
and retaining contraction partners, followed by a subgroup mapping table.}

\paragraph{1. Why reassociation requires splitting groups.}
Dot reassociation changes $(XY)Z$ to $X(YZ)$. For generalized dots it must also
specify the axes contracted by each replacement. Consider this rule, using
\code{dot} for a dialect's generalized dot:
\begin{codeblock}[H]
\begin{lstlisting}[language=TEPL]
rule associate {
  (dot[@outer] (dot[@inner] X Y) Z)
  => (dot[auto] X (dot[auto] Y Z))
  where {
    @outer.lhs_batching_dimensions == @inner->B
  }
}
\end{lstlisting}
\end{codeblock}
Incompatible batch routing can change a reassociation's result. The guard
requires the outer dot's left batch list to equal the positions of the inner
dot's batch group in its result, in the same order. The expression
\code{@inner->B} denotes those positions. This restricts applicability;
the completed replacement still requires the checks in
Sections~\ref{sec:synthesis} and~\ref{sec:correctness}.

Let $T=\operatorname{dot\_general}(X,Y)$. The inner dot contracts $K$ and
yields $P_T=B\catseq I\catseq J$, with free groups $I$ from \code{X} and
$J$ from \code{Y}. Write $c_L^{\mathrm{outer}},c_R^{\mathrm{outer}}$ for the
outer contraction lists in \code{T} and \code{Z}. Their shared group is
\begin{equation}
 L=\gath(P_T,c_L^{\mathrm{outer}}).
\end{equation}
Since $B$ remains batch, the left contraction selects only from $I,J$.
Unlike addition's single group, this producer has groups from different
operands. The outer dot can select part of each, mixing the selected axes
in $L$. To build \code{dot_general(Y,Z)}, the compiler must separate the
portion from $J$, which can be contracted there, from the partners of $I$,
which must survive for the final dot with \code{X}.

\paragraph{2. Find the selected part of each producer group.}
For each producer group, retain both within-group positions and contraction
slots. Group $I$ occupies
\begin{equation}
 p_I=|B|+\rho(I)
\end{equation}
in $P_T$. Compare this list to the outer selection using
\begin{equation}
 (u_I,v_I)=\mpo(c_L^{\mathrm{outer}},p_I).
\end{equation}
For distinct-entry sequences $a,b$, the operation
$\mpo(a,b)=(u,v)$ scans $a$ in order. Whenever $a_i=b_j$, it appends $i$ to
$u$ and $j$ to $v$. It returns positions in the input lists, rather than
the shared values. Hence $\gath(a,u)=\gath(b,v)$, with both results empty
when there is no overlap. For example,
\begin{equation}
\begin{aligned}
 a&=\langle3,0,4\rangle,&b&=\langle0,2,3\rangle,\\
 u&=\langle0,1\rangle,&v&=\langle2,0\rangle.
\end{aligned}
\end{equation}
The two gathers produce $\langle3,0\rangle$, retaining $a$'s order.
For symbolic inputs, the results remain position expressions until they can
be evaluated or simplified.

Here $u_I$ locates the selected axes in the contraction list, and $v_I$
locates them within $I$:
\begin{equation}
 \gath(c_L^{\mathrm{outer}},u_I)=\gath(p_I,v_I).
\end{equation}
The selected subgroup is therefore
\begin{equation}
 \Isel=\gath(I,v_I)=\gath(L,u_I).
\end{equation}
The remaining view uses an ordered complement:
\begin{equation}
 w_I=\ocomp(\rho(I),v_I),\qquad \Irem=\gath(I,w_I).
\end{equation}
The selected view follows outer contraction order, whereas the remaining
view follows $I$'s original order. Their concatenation need not reproduce
$I$, so the position sequences must be retained. The same construction
applies to $J$ using $p_J=|B|+|I|+\rho(J)$. For either $G\in\{I,J\}$,
$u_G,v_G$ locate its selected portion, and
$w_G=\ocomp(\rho(G),v_G)$ locates its remaining portion. The lengths and
selected portions remain symbolic.

Partitioning adds the following relationships:
\begin{center}
\small
\begin{tabular}{@{}lcccc@{}}\toprule
Group & In $P_T$ & In $I$ & In $J$ & In $L$\\\midrule
$\Isel$ & $\gath(p_I,v_I)$ & $v_I$ & -- & $u_I$\\
$\Irem$ & $\gath(p_I,w_I)$ & $w_I$ & -- & --\\
$\Jsel$ & $\gath(p_J,v_J)$ & -- & $v_J$ & $u_J$\\
$\Jrem$ & $\gath(p_J,w_J)$ & -- & $w_J$ & --\\\bottomrule
\end{tabular}
\end{center}
The $L$ column supplies the connection to \code{Z} through contraction slots.

\paragraph{3. Map selected axes to \code{Z}.}
The outer dot pairs its left and right lists entry by entry. Selecting slots
$u_I$ and $u_J$ from its right list therefore gives the positions in \code{Z}
paired with the selected subgroups:
\begin{equation}
 z_I=\gath(c_R^{\mathrm{outer}},u_I),\qquad
 z_J=\gath(c_R^{\mathrm{outer}},u_J).
\label{eq:z-partners}
\end{equation}
In Figure~\ref{fig:axis-partitioning}, \code{Z} has axes $[z_0,f,z_1]$ and
right contraction positions $[2,0]$. The left list selects $[j_1,i_0]$.
Axis $j_1$ occupies contraction slot $0$ and pairs with $z_1$ at position $2$;
$i_0$ occupies slot $1$ and pairs with $z_0$ at position $0$.
The selected groups thus retain all three descriptions:
\begin{center}
\begin{tabular}{@{}lccc@{}}\toprule
Selected group & In producer group & Slots in $L$ & In $Z$\\\midrule
$\Isel$ & $v_I$ & $u_I$ & $z_I$\\
$\Jsel$ & $v_J$ & $u_J$ & $z_J$\\\bottomrule
\end{tabular}
\end{center}
Combining these entries with producer mappings describes the original
computation in terms of the refined groups without losing operand pairings.

\paragraph{4. Extend the original mapping table.}
The original graph is $T=\operatorname{dot\_general}(X,Y)$ followed by
$O=\operatorname{dot\_general}(T,Z)$. Write $b_X,k_X,i_X$ for the known
positions of $B,K,I$ in \code{X}, and $b_Y,k_Y,j_Y$ for those of $B,K,J$ in
\code{Y}. Composition locates each subgroup: for example,
$\gath(i_X,v_I)$ locates $\Isel$ in \code{X}. In \code{Z}, let $b_Z$ be the
batch positions and $f_Z$ the positions of free group $F$.

The inner reduction removes $K$ from $T$; the outer reduction removes
$\Isel,\Jsel$ from $O$. The surviving order is
\begin{equation}
 P_O=B\catseq\Irem\catseq\Jrem\catseq F.
\label{eq:source-output}
\end{equation}
Table~\ref{tab:refined-mappings} combines the refined relationships. To keep
entries readable, define
\begin{equation}
\begin{aligned}
 i_s&=\gath(i_X,v_I),&i_r&=\gath(i_X,w_I),\\
 j_s&=\gath(j_Y,v_J),&j_r&=\gath(j_Y,w_J),\\
 t_{Is}&=\gath(p_I,v_I),&t_{Ir}&=\gath(p_I,w_I),\\
 t_{Js}&=\gath(p_J,v_J),&t_{Jr}&=\gath(p_J,w_J),\\
 o_I&=|B|+\rho(\Irem),\\
 o_J&=|B|+|\Irem|+\rho(\Jrem),\\
 o_F&=|B|+|\Irem|+|\Jrem|+\rho(F).
\end{aligned}
\end{equation}
These are abbreviations for the position expressions, not additional
attributes to be supplied by the user.
\begin{table}[htbp]
\caption{Refined mappings across the original reassociation graph. Each entry
gives positions in the column's interface. A dash indicates no mapping in
this example.}
\label{tab:refined-mappings}
\centering
\begin{tabular}{@{}lccccc@{}}\toprule
Group/interface & $X$ & $Y$ & $T$ & $Z$ & $O$\\\midrule
Axis interface & $P_X$ & $P_Y$ & $P_T$ & $P_Z$ & $P_O$\\
$B$ & $b_X$ & $b_Y$ & $\rho(B)$ & $b_Z$ & $\rho(B)$\\
$K$ & $k_X$ & $k_Y$ & -- & -- & --\\
$\Isel$ & $i_s$ & -- & $t_{Is}$ & $z_I$ & --\\
$\Irem$ & $i_r$ & -- & $t_{Ir}$ & -- & $o_I$\\
$\Jsel$ & -- & $j_s$ & $t_{Js}$ & $z_J$ & --\\
$\Jrem$ & -- & $j_r$ & $t_{Jr}$ & -- & $o_J$\\
$F$ & -- & -- & -- & $f_Z$ & $o_F$\\\bottomrule
\end{tabular}
\end{table}

The procedure generalizes to other producer--consumer pairs: compare position
sequences on the shared interface, identify overlap, and add mappings from
each overlapping portion into both parent groups. Each partition is internally
disjoint, while groups in different partitions can overlap.

\FloatBarrier
\subsubsection{General Construction Algorithm}
\label{sec:mapping-algorithm}
The construction algorithm processes producers before consumers, connects each
operation's bindings to its inputs, and extends mappings to the output.
Overlap introduces the finer subgroup relationships described above.
The input is a source graph and its operation definitions; the output comprises
tensor interfaces, composed value expressions, established mappings,
and partitions.

\paragraph{State.}
The state is $\Sigma=(V,M,\Pi)$. For each tensor $t$,
$V[t]=(P_t,e_t)$ stores its ordered interface and value expression.
$M$ contains individual mapping entries $(P,G,p)$. The partition collection
$\Pi$ records which groups cover an interface completely and without overlap:
\begin{equation}
 \pi=\bigl(P,\langle(G_1,p_1),\ldots,(G_m,p_m)\rangle\bigr).
\end{equation}
An interface can have several partitions because different operations group
its axes differently. A partition changes the grouping used in reasoning;
it does not change tensor shape, rank, or axis order. $M$ locates the groups,
whereas $\Pi$ says which groups jointly cover the interface exactly once.

For example, if addition produces $T$ with $S=\langle i,j\rangle$, then
$V[T]=(P_T,\operatorname{Add}(X,Y))$ with $P_T=S$,
$M$ includes $(P_T,S,\langle0,1\rangle)$, and $\Pi$ includes the single-member
partition $(P_T,\langle(S,\langle0,1\rangle)\rangle)$. A consumer may split
$S$ into $I=\langle i\rangle$ and $J=\langle j\rangle$. Then $M$ gains their
position mappings and $\Pi$ gains the two-member partition alongside the
original one. A shape $\langle d_i,d_j\rangle$ remains unchanged.

\paragraph{Invariants.}
Every mapping must select its group in the stated order:
\begin{equation}
 (P,G,p)\in M\ \Longrightarrow\ \gath(P,p)=G.
 \label{eq:mapping-invariant}
\end{equation}
Every partition member satisfies the same invariant, and
$p_1\catseq\cdots\catseq p_m$ must be a permutation of $\rho(P)$.
Thus each axis occurs exactly once within one partition, although different
partitions may overlap. In the addition example,
$\gath(\langle i,j\rangle,\langle0,1\rangle)=S$, and $\langle0,1\rangle$
covers every position. The sequence $\langle0,0\rangle$ would repeat one
axis and omit the other, violating the partition invariant.

\begin{paperalgorithm}[H]
\caption{Build symbolic axis mapping.}\label{alg:mapping}
\begin{lstlisting}[language=Python]
def build_symbolic_axis_mapping(graph):
    V = initialize_inputs(graph)
    M = set()       # Entries (interface, group, positions)
    partitions = {}

    for op in producer_first(graph):
        inputs = [V[t] for t in op.inputs]
        groups, bindings = resolve_bindings(op, inputs)
        aligned = []

        for (P, expression), operand_bindings in zip(inputs, bindings):
            Q = build_interface(P, operand_bindings)
            for group, positions in operand_bindings:
                M.add((Q, group, positions))
            refine_and_register(Q, operand_bindings, M, partitions)
            expression = transport(expression, P, Q, M, partitions)
            aligned.append((Q, expression))

        Pout = evaluate(op.output_axes, groups)
        output_bindings = list(output_positions(op, groups))
        for group, positions in output_bindings:
            M.add((Pout, group, positions))
        refine_and_register(Pout, output_bindings, M, partitions)
        V[op] = (Pout, compose(op.value, aligned, groups))

    return V, M, partitions
\end{lstlisting}
\end{paperalgorithm}

\code{initialize_inputs} gives each graph input an interface and read
expression. \code{producer_first} ensures descriptions are available before
use. \code{resolve_bindings} evaluates symbolic position expressions and
establishes validity and pairing conditions. It reuses producer coordinates
where applicable and otherwise introduces local groups. Each operand's
bindings $(G_i,p_i)$ cover its axes.

\code{build_interface} reuses $P$ when the required groups already occupy the
required positions. Otherwise it constructs
\begin{equation}
 Q=\gath\bigl(G_1\catseq\cdots\catseq G_m,
        \operatorname{inverse}(p_1\catseq\cdots\catseq p_m)\bigr).
\end{equation}
The inverse arranges collected groups in tensor dimension order.
\code{refine_and_register} registers the new partition and compares it
with earlier partitions on the same interface. It uses \code{matchPositions}
to add overlap mappings. \code{transport} updates the producer's read
coordinates from $P$ to $Q$ and carries mappings and partitions into that
description for the current consumer. Equal interfaces need no coordinate
update. Finally, \code{evaluate} and \code{output_positions} derive the
output interface and positions from the declared order, and \code{compose}
combines transported expressions using the operation's value definition.

\paragraph{Deriving mappings on demand.}
A missing mapping can be recovered by following established relationships:
\begin{equation}
\begin{aligned}
 (P,G,p)\in M\ \land\ (G,H,q)\in M
 \quad\Longrightarrow\quad
 \gath(P,\gath(p,q))=H.
\end{aligned}
\end{equation}
Lookup can return $\gath(p,q)$ or register the corresponding entry. Some
relationships propagate during construction and others are derived when
queried; enumerating every possible mapping in advance is unnecessary.
Local bindings and output order establish local mappings, paired position
matches establish overlaps, and Equation~\eqref{eq:gather-composition}
composes them. These operations are designed to maintain the invariant;
they do not dispense with checking the completed rewrite.
